\documentclass[10pt,a4paper]{amsart}
\usepackage[margin=19mm]{geometry}
\usepackage{amsmath,amssymb,mathtools}
\usepackage[scr=rsfs,cal=euler]{mathalfa}
\usepackage{xfrac}
\usepackage[unicode,hidelinks,bookmarks=false]{hyperref}
\newcommand{\ie}{i.e.\ }
\newcommand{\cf}{cf.\ }
\newcommand{\resp}{respectively\ }
\newcommand{\Subsection}{\subsection}
\providecommand{\nobreakdash}{\nobreak\hbox{-}}
\setlength{\emergencystretch}{2em}
\multlinegap0pt
\usepackage{bm}
\usepackage{bbm}
\usepackage{dsfont}
\usepackage{xspace}
\usepackage{cancel}
\usepackage{mathtools}
\usepackage{amsthm}
\makeatletter
\@ifundefined{claim}{\newtheorem{claim}{Claim}}{}
\@ifundefined{corollary}{\newtheorem{corollary}{Corollary}}{}
\@ifundefined{lemma}{\newtheorem{lemma}{Lemma}}{}
\makeatother
\title[An arithmetic algebraic regularity lemma]{An arithmetic algebraic regularity lemma}
\author{Anand Pillay, Atticus Stonestrom}
\date{Source: arXiv:2412.11206v2 (CC BY 4.0)}
\begin{document}
\maketitle

\noindent\emph{Source and licence.} An arithmetic algebraic regularity lemma. Version arXiv:2412.11206v2 (CC BY 4.0), distributed under the Creative Commons Attribution 4.0 International licence, as recorded in the arXiv licence metadata for this exact version and shown on the abstract page https://arxiv.org/abs/2412.11206v2. Full rights notice: licenses/arxiv-2412.11206-CC-BY-4.0.txt.\par\smallskip

\noindent\textit{Editorial scope.} Counted unit: the Lemma whose source label is \texttt{algebraic\_regularity\_lemma} in the article (source0.tex lines 369--371, source label \texttt{algebraic\_regularity\_lemma}), delivered together with the article's own complete original proof of it (source0.tex lines 372--407, one \texttt{proof} environment of 36 lines). No mathematics is written, repaired or re-derived: statement and proof are the article's own text line for line. Required context retained uncounted: HF\_finitary (lines 341--347). Every definition, display and label of those lines is retained unchanged. Omitted from this package and not redistributed: 1--340, 348--368, 408--656. Nothing inside a retained excerpt is omitted or altered: every retained body is delivered exactly as the article's lines are written, so no omission receipt of any kind accompanies this package. Citations: the retained excerpts carry no citation command at all, so no bibliography entry of the article is transcribed and no reference is redistributed. Reference adaptation: none was needed and none was made; every reference inside the delivered unit resolves inside it, so no unresolved reference or citation is produced. Printed slips kept literal: no printed word, symbol, hypothesis or number of the counted statement or of its proof was altered. Layout and class: the article's own class is \texttt{article} with packages including preamble, courier, natbib; the wrapper is the lane's pinned \texttt{amsart} editorial wrapper at 10pt on A4, which restates the theorem-like environments the retained text needs and adds the article's own macro definitions that the retained text needs (none), with extra wrapper packages (bm, bbm, dsfont, xspace, cancel, mathtools). The delivered numbering is the unit's own. Assets: no retained span contains a figure environment, a graphics-inclusion command, a PDF-page inclusion, a table or a TikZ picture; no image file is copied into this package, the package declares no asset dependency and no image file is redistributed with it. Licence: version arXiv:2412.11206v2 of this article is distributed under the Creative Commons Attribution 4.0 International licence, as recorded in the arXiv licence metadata for this exact version and shown on the abstract page https://arxiv.org/abs/2412.11206v2. A full rights notice is delivered at licenses/arxiv-2412.11206-CC-BY-4.0.txt. Characters: every retained excerpt is pure ASCII, so the delivered engine, XeLaTeX with its default Latin Modern text font, reports no missing character.
\begin{corollary}\label{HF_finitary}
    Fix $M>0$. Then there is $C=C(M)>0$ such that the following holds for any finite field $\mathbf{F}$. If $G$ is a definable group in $\mathbf{F}$ of complexity $\leqslant M$, and $D\subseteq G$ is a definable subset of $G$ of complexity $\leqslant M$, then there is a definable normal subgroup $H\leqslant G$, of complexity at most $C$, with $[G:H]\leqslant C$ and with the following properties:
    \begin{enumerate}
        \item For any cosets $V,W$ of $H$, there is a rational number $r\leqslant 1$ such that, for all but $\leqslant C|\mathbf{F}|^{-1/2}|G|^2$ many pairs $(b,c)\in V\times W$, $\big||b^{-1}D\cap c^{-1}D|-r|G|\big|\leqslant C|\mathbf{F}|^{-1/2}|G|$.
        \item For any cosets $V,W$ of $H$, there is a rational number $r\leqslant 1$ such that, for all but $\leqslant C|\mathbf{F}|^{-1/2}|G|^2$ many pairs $(b,c)\in V\times W$, $\big||Db\cap Dc|-r|G|\big|\leqslant C|\mathbf{F}|^{-1/2}|G|$.
    \end{enumerate}
 \end{corollary}

\begin{lemma}\label{algebraic_regularity_lemma}
    Fix any $M>0$. Then there is $C=C(M)>0$ such that the following holds for any finite field $\mathbf{F}$. If $G$ is a definable group in $\mathbf{F}$ of complexity $\leqslant M$, and $D\subseteq G$ is a definable subset of $G$ of complexity $\leqslant M$, then there is a normal subgroup $H\leqslant G$ with $[G:H]\leqslant C$ and such that the following holds. For any cosets $V,W$ of $H$, the graph $(V,W,xy^{-1}\in D)$ is weakly $C|\mathbf{F}|^{-1/4}$-regular; ie, for any $A\subseteq V,B\subseteq W$, the value $|\{(a,b)\in A\times B:ab^{-1}\in D\}|$ differs from $$|\{(v,w)\in V\times W:vw^{-1}\in D\}|\frac{|A||B|}{|V||W|}$$ by at most $C|\mathbf{F}|^{-1/4}|H|^2$.
\end{lemma}

\begin{proof}
    Fix $M>0$ and let $C_1$ be given by Corollary \ref{HF_finitary}. Suppose that $\mathbf{F}$ is a finite field, that $G$ is a definable group in $\mathbf{F}$ of complexity $\leqslant M$, and that $D\subseteq G$ is a definable set of complexity $\leqslant M$. Let $H$ be given by Corollary \ref{HF_finitary}.

\begin{claim}\label{mean_0_individual}
    There is a constant $C_2$ depending only on $M$ such that, for any $U$ a coset of $H$, and any $f:U\to\mathbb{R}$ of mean $0$ and with $||f||_\infty\leqslant 1$,
    \begin{align*}
        &\sum_{g\in G}\left(\sum_{u\in U}1_{[ug\in D]}f(u)\right){}^2\text{ and }\sum_{g\in G}\left(\sum_{u\in U}1_{[gu^{-1}\in D]}f(u)\right){}^2
    \end{align*}are each bounded above by $C_2|\mathbf{F}|^{-1/2}|H|^3$.
\end{claim}
\begin{proof}
    We show the first bound, as the second bound is proved symmetrically using condition (2) of Corollary \ref{HF_finitary}; let $\Delta$ be the quantity we wish to bound. Expanding out the square in the sum gives $$\Delta=\sum_{g\in G}\sum_{u,v\in U}1_{[ug\in D]}1_{[vg\in D]}f(u)f(v),$$ which is just $\sum_{(u,v)\in U\times U}|u^{-1}D\cap v^{-1}D|f(u)f(v)$. By condition (1) of Corollary \ref{HF_finitary}, there is some rational $r\leqslant 1$ such that the set $$F:=\{(u,v)\in U\times U:\left||u^{-1}D\cap v^{-1}D|-r|G|\right|>C_1|\mathbf{F}|^{-1/2}|G|\}$$ has size $\leqslant C_1|\mathbf{F}|^{-1/2}|G|^2$. Now, since $||f||_\infty\leqslant 1$, $\big|r|G|f(u)f(v)-|u^{-1}D\cap v^{-1}D|f(u)f(v)\big|$ is bounded above by $C_1|\mathbf{F}|^{-1/2}|G|$ for all $(u,v)\in (U\times U)\setminus F$, whence $$\bigg|\sum_{(u,v)\in (U\times U)\setminus F}r|G|f(u)f(v)-\sum_{(u,v)\in (U\times U)\setminus F}|u^{-1}D\cap v^{-1}D|f(u)f(v)\bigg|$$ is bounded above by $|(U\times U)\setminus F|C_1|\mathbf{F}|^{-1/2}|G|\leqslant C_1|\mathbf{F}|^{-1/2}|G|^3\leqslant C_1^4|\mathbf{F}|^{-1/2}|H|^3$, where in the last inequality we used that $|G|\leqslant C_1|H|$.
    
    Also, since $D\subseteq G$ and $r\leqslant 1$, we have $\left||u^{-1}D\cap v^{-1}D|-r|G|\right|\leqslant |G|$ for all $(u,v)\in G\times G$, and in particular for all $(u,v)\in F$, so that $$\bigg|\sum_{(u,v)\in F}r|G|f(u)f(v)-\sum_{(u,v)\in  F}|u^{-1}D\cap v^{-1}D|f(u)f(v)\bigg|$$ is bounded above by $|F||G|\leqslant C_1|\mathbf{F}|^{-1/2}|G|^2|G|\leqslant C_1^4|\mathbf{F}|^{-1/2}|H|^3$, where in the last inequality we again used that $|G|\leqslant C_1|H|$.
    
    Taking $C_2=2C_1^4$ and applying the triangle inequality to the two estimates above gives that $\Delta$ differs from $\sum_{(u,v)\in U\times U}r|G|f(u)f(v)$ by at most $C_2|\mathbf{F}|^{-1/2}|H|^3$. But $f$ has mean $0$ so the latter sum is $0$ and the claim follows.
\end{proof}

\begin{claim}\label{mean_0_both}
    There is a constant $C_3$ depending only on $M$ such that, for any cosets $V,W$ of $H$, and any $f:V\to\mathbb{R}$ and $g:W\to\mathbb{R}$ with $||f||_\infty,||g||_\infty\leqslant 1$, if at least one of $f,g$ has mean $0$ then $|\sum_{(v,w)\in V\times W}1_{[vw^{-1}\in D]}f(v)g(w)|\leqslant C_3|\mathbf{F}|^{-1/4}|H|^2$.
\end{claim}

\begin{proof}
    Suppose for instance that $f$ has mean $0$; the other case is symmetric using the other half of Claim \ref{mean_0_individual}. We have
    \begin{align*}
    \left(\sum_{(v,w)\in V\times W}f(v)g(w)1_{[vw^{-1}\in D]}\right){}^2&=\left(\sum_{w\in W}\left(\sum_{v\in V}f(v)g(w)1_{[vw^{-1}\in D]}\right)\right){}^2\\
    &\leqslant |W|\sum_{w\in W}\left(\sum_{v\in V}f(v)g(w)1_{[vw^{-1}\in D]}\right){}^2 \\
    &=|W|\sum_{w\in W}g(w)^2\left(\sum_{v\in V}f(v)1_{[vw^{-1}\in D]}\right){}^2 \\
    &\leqslant |W|\sum_{w\in W}\left(\sum_{v\in V}f(v)1_{[vw^{-1}\in D]}\right){}^2 \\
    &\leqslant |W|\sum_{w\in G}\left(\sum_{v\in V}f(v)1_{[vw^{-1}\in D]}\right){}^2
\end{align*} where the first inequality follows from Cauchy-Schwarz, and the second follows from the fact that $||g||_\infty\leqslant 1$. But now $f$ has mean $0$, so that by the first bound in Claim \ref{mean_0_individual} the left hand side is bounded above by $|W|C_2|\mathbf{F}|^{-1/2}|H|^3= C_2|\mathbf{F}|^{-1/2}|H|^4$. So taking $C_3=\sqrt{C_2}$ works.
\end{proof}

Let $C_3$ be given by Claim 2. Suppose that $V,W$ are cosets of $H$ and that $A\subseteq V$ and $B\subseteq W$. Define $f:V\to \mathbb{R}$ and $g:W\to\mathbb{R}$ by $f(v)=1_{A}(v)-|A|/|H|$ and $g(w)=1_{B}(w)-|B|/|H|$. By Claim 2, $\sum_{(v,w)\in V\times W}\frac{|A|}{|H|}g(w)1_{[vw^{-1}\in D]}$ and $\sum_{(v,w)\in V\times W}f(v)\frac{|B|}{|H|}1_{[vw^{-1}\in D]}$ and $\sum_{(v,w)\in V\times W}f(v)g(w)1_{[vw^{-1}\in D]}$ are all bounded in magnitude by $C_3|\mathbf{F}^{-1/4}||H|^2$. Also clearly $\sum_{(v,w)\in V\times W}\frac{|A|}{|H|}\frac{|B|}{|H|}1_{[vw^{-1}\in D]}=|\{(v,w)\in V\times W:vw^{-1}\in D\}|\frac{|A||B|}{|V||W|}$. But the sum of those four terms is 
\begin{align*}\sum_{(v,w)\in V\times W}\left(f(v)+\frac{|A|}{|H|}\right)\left(g(w)+\frac{|B|}{|H|}\right)1_{[vw^{-1}\in D]}&=\sum_{(v,w)\in V\times W}1_A(v)1_B(w)1_{[vw^{-1}\in D]}\\&=|\{(a,b)\in A\times B:ab\in D\}|.
\end{align*}Thus, applying the triangle inequality to the sum of the three terms bounded above, $|\{(a,b)\in A\times B:ab\in D\}|$ differs from $|\{(v,w)\in V\times W:vw^{-1}\in D\}|\frac{|A||B|}{|V||W|}$ by at most $3C_3|\mathbf{F}|^{-1/4}|H|^2$, and so taking $C=3C_3$ gives the result.
\end{proof}

\end{document}
